\documentclass[border=10pt,crop=true]{standalone}
\usepackage{circuitikz}
\usetikzlibrary{calc}
% Shared style for 电子学一 Chapter 2 op-amp circuit figures (CircuiTikZ).
\usetikzlibrary{calc}

\ctikzset{
  bipoles/length=0.95cm,
  bipoles/thickness=1,
  resistors/scale=1,
  capacitors/scale=1,
  label distance=2.5mm,
}

\tikzset{
  opfig/.style={font=\small},
  opfigThin/.style={font=\scriptsize},
}

% Geometry (cm) — keep identical across all op-amp schematics
\def\OpInLen{1.55}    % label → input terminal
\def\OpInGap{0.08}    % terminal → wire to op amp +
\def\OpAmpWire{1.00}  % terminal side → op amp + anchor
\def\OpOutLen{1.05}   % op amp out → output terminal
\def\OpJuncL{0.50}    % v_- → feedback junction (left)
\def\OpFbDown{0.55}   % v_- straight down before R_1
\def\OpRvert{1.20}    % R_1 vertical span to ground
\def\OpInSep{0.55}    % dual-input spacing (v_1 / v_2)
\def\OpFbStub{0.40}   % follower: stub right of output before dropping
\def\OpFbClear{0.22}  % follower: below op-amp south for horizontal feedback


\ctikzset{resistors/scale=0.8, label distance=3mm}

% 布局参数（cm）
\def\OpIaXa{2.45}
\def\OpIaYpair{1.25}
\def\OpIaXin{-1.22}
\def\OpIaXinv{1.50}
\def\OpIaRtwoSpan{0.68}
\def\OpIaXaThree{7.05}
\def\OpIaRthreeLen{2.05}
\def\OpIaRfourDrop{1.25}
\def\OpIaYbus{-2.40}
\def\OpIaXfbR{8.05}

\begin{document}
\begin{circuitikz}[american, scale=1.0, transform shape]

  % ==================== 1. 三个运算放大器 ====================
  \node[op amp, noinv input up](OA1) at (\OpIaXa, \OpIaYpair) {};
  \node[op amp, noinv input down](OA2) at (\OpIaXa, -\OpIaYpair) {};
  \node[op amp, noinv input up](OA3) at (\OpIaXaThree, 0) {};

  % ==================== 2. 左上 OA1 ====================
  \path let \p1=(OA1.+), \p2=(OA1.-), \p3=(OA1.out) in
    coordinate (V1t) at (\OpIaXin, \y1)
    coordinate (N1) at (\OpIaXinv, \y2)
    coordinate (Out1) at (\x3, \y3);

  \draw (OA1.+) -- (V1t) to[short, o-] ++(-\OpInGap, 0) node[left, opfig]{$v_1$};
  \draw (OA1.-) -| (N1);
  \node[circ, fill=black, inner sep=0.9pt] at (N1) {};

  \coordinate (R2h1) at (\OpIaXinv, \OpIaYpair-\OpIaRtwoSpan);
  \draw (N1) |- (R2h1)
    to[R, l^=$R_2$] (R2h1 -| Out1) |- (Out1);

  % ==================== 3. 左下 OA2 ====================
  \path let \p1=(OA2.+), \p2=(OA2.-), \p3=(OA2.out) in
    coordinate (V2t) at (\OpIaXin, \y1)
    coordinate (N2) at (\OpIaXinv, \y2)
    coordinate (Out2) at (\x3, \y3);

  \draw (OA2.+) -- (V2t) to[short, o-] ++(-\OpInGap, 0) node[left, opfig]{$v_2$};
  \draw (OA2.-) -| (N2);
  \node[circ, fill=black, inner sep=0.9pt] at (N2) {};

  \coordinate (R2h2) at (\OpIaXinv, -\OpIaYpair+\OpIaRtwoSpan);
  \draw (N2) |- (R2h2)
    to[R, l_=$R_2$] (R2h2 -| Out2) |- (Out2);

  % ==================== 4. 输入级 R_1（竖向，标签在左侧）====================
  \draw (N1) to[R, l_=$R_1$] (N2);

  % ==================== 5. 末级差动放大器 ====================
  \path let \p1=(OA3.+), \p2=(OA3.-), \p3=(OA3.out) in
    coordinate (PlusJ) at ($(\x1, \y1)+(-\OpIaRthreeLen, 0)$)
    coordinate (MinusJ) at ($(\x2, \y2)+(-\OpIaRthreeLen, 0)$)
    coordinate (FbL) at (\x2, \OpIaYbus)
    coordinate (FbR) at (\OpIaXfbR, \OpIaYbus);

  % --- 上路：Out1 → R_3 → OA3.+ ；从 (+) 垂直接 R_4 → 地（无中间叠层）---
  \draw (Out1) -| (PlusJ);
  \draw (PlusJ) to[R, l=$R_3$] (OA3.+);
  \draw (OA3.+) to[R, l_=$R_4$] ++(0, -\OpIaRfourDrop)
    node[ground, scale=0.85]{};

  % --- 下路：Out2 → R_3 → OA3.-（标签在电阻下方）---
  \draw (Out2) -| (MinusJ);
  \draw (MinusJ) to[R, l_=$R_3$] (OA3.-);

  % --- 负反馈：(-) 直下母线 → 水平 R_4（标签在上）→ 升至输出 ---
  \draw (OA3.-) -- (FbL);
  \draw (FbL) to[R, l^=$R_4$] (FbR);
  \draw (FbR) |- (OA3.out);

  % ==================== 6. 输出 v_o ====================
  \draw (OA3.out) to[short, -o] ++(\OpOutLen, 0) node[right, opfig]{$v_o$};

\end{circuitikz}
\end{document}
